\subsection{Algebra of limits}

When individual limits exist, sums, differences, products, and quotients can often be handled term by term.

\begin{theorembox}[Elementary limit laws]
If
\[
\lim_{x\to a}f(x)=L,
\qquad
\lim_{x\to a}g(x)=M,
\]
then
\[
\lim_{x\to a}(f+g)=L+M,
\qquad
\lim_{x\to a}(fg)=LM,
\]
and, when \(M\ne0\),
\[
\lim_{x\to a}\frac{f(x)}{g(x)}=\frac{L}{M}.
\]
The same principles hold for sequence limits.
\end{theorembox}

\begin{examplebox}[Using laws before using tricks]
For
\[
\lim_{x\to2}\frac{x^2+3x}{x+1},
\]
the denominator tends to \(3\ne0\), so direct use of the laws gives
\[
\frac{2^2+3\cdot2}{2+1}=\frac{10}{3}.
\]
No transformation is needed.
\end{examplebox}

\begin{remarkbox}
The expressions \(0/0\), \(\infty/\infty\), and \(\infty-\infty\) are not results. They indicate that the elementary laws cannot be applied directly and that the expression must be reorganized.
\end{remarkbox}
